\documentclass[12pt,letterpaper]{article}     % Tipo de documento y otras especificaciones
\usepackage[utf8]{inputenc}                   % Para escribir tildes y eñes
\usepackage[spanish]{babel}                   % Para que los títulos de figuras, tablas y otros estén en español
\addto\captionsspanish{\renewcommand{\tablename}{Tabla}}					% Cambiar nombre a tablas
\addto\captionsspanish{\renewcommand{\listtablename}{Índice de tablas}}		% Cambiar nombre a lista de tablas
\usepackage{geometry}                         
\geometry{left=18mm,right=18mm,top=21mm,bottom=21mm} % Tamaño del área de escritura de la página
\usepackage{ucs}
\usepackage{amsmath}      % Los paquetes ams son desarrollados por la American Mathematical Society
\usepackage{amsfonts}     % y mejoran la escritura de fórmulas y símbolos matemáticos.
\usepackage{amssymb}
\usepackage{graphicx}     % Para insertar gráficas
\usepackage[lofdepth,lotdepth]{subfig}	% Para colocar varias figuras
\usepackage{unitsdef}	  % Para la presentación correcta de unidades
\usepackage{pdfpages}   %incluir paginas de pdf externo, para los anexos
\usepackage{appendix}   %para los anexos
\renewcommand{\unitvaluesep}{\hspace*{4pt}}	% Redimensionamiento del espacio entre magnitud y unidad
\usepackage[colorlinks=true,urlcolor=blue,linkcolor=black,citecolor=black]{hyperref}     % Para insertar hipervínculos y marcadores
\usepackage{float}		% Para ubicar las tablas y figuras justo después del texto
\usepackage{booktabs}	% Para hacer tablas más estilizadas
\batchmode
\bibliographystyle{plain} 
\pagestyle{plain} 
\pagenumbering{arabic}
\usepackage{lastpage}
\usepackage{fancyhdr}	% Para manejar los encabezados y pies de página
\pagestyle{fancy}		% Contenido de los encabezados y pies de pagina


\lhead{IE0308 - Laboratorio Eléctrico I}
\chead{}
\rhead{Experimento }	% Aquí va el numero de experimento, al igual que en el titulo
\lfoot{Escuela de Ingeniería Eléctrica}
\cfoot{\thepage\ de \pageref{LastPage}}
\rfoot{Universidad de Costa Rica}


\author{Nombre1, Carné1 \\ Nombre2, Carné2 \\ {\small Grupo 02}\\ Profesor:  \vspace*{3.0in}}
\title{Universidad de Costa Rica\\{\small Facultad de Ingeniería\\Escuela de Ingeniería Eléctrica\\IE0308 – Laboratorio Eléctrico I\\II ciclo 2013\\\vspace*{0.55in} Anteproyecto}\\ Experimento \vspace*{1.35in}}
%\date{}  				

%%%%%%%%%%%%%%%%
\begin{document}	% Inicio del documento
%%%%%%%%%%%%%%%%

\pdfbookmark[1]{Portada}{portada} 	% Marcador para el título

\maketitle							% Título


\tableofcontents
\newpage
\listoffigures
\listoftables

\newpage

%%%%%%%%%%%%%%%%%%%
\section{Objetivos}
%%%%%%%%%%%%%%%%%%%

\subsection{Objetivo general}
%%%%%%%%%%%%%%%%%%%%%%%%%%%%%

Objetivo descrito en la Guía del Experimento.


\subsection{Objetivos específicos}
%%%%%%%%%%%%%%%%%%%%%%%%%%%%%%%%%%

\begin{itemize}
\item Objetivos específicos propuestos por los estudiantes.
\item Son exhaustivos y basados en la nota teórica, el procedimiento y los resultados esperados del experimento.
\end{itemize}

\newpage
%%%%%%%%%%%%%%%%%%%%%%
\section{Nota teórica}
%%%%%%%%%%%%%%%%%%%%%%

\setcounter{secnumdepth}{4}



\subsection{Titulo.}

\subsubsection{Subtitulo.}

Debe incluir todas las justificaciones teóricas que el diseño propuesto amerite, además de las investigaciones propuestas en el enunciado de cada experimento. Se debe hacer referencia a las citas y figuras incluidas, según las fuentes bibliográficas utilizadas. 

Las imágenes deben hacerse referencia como Figura 1 o Figura 2, etc. Ejemplo: “Tal y como se muestra en la Figura 1, hay una relación lineal entre la tensión y la corriente eléctrica”. No deben hacerse referencias del tipo “En la siguiente figura” o “Como se vio en la figura anterior” porque esto puede ser un problema si se cambia el orden del documento o la posición de las figuras. Para esto se utiliza el label de la imagen con el comando ref. Además se hace una referencia bibliografica con el comando cite.

\begin{figure}[H]
\centering
\includegraphics[width=0.3\textwidth]{./imagenes/fig1.png} 
\caption{Circuito de Prueba} \cite{Boylestad}
\label{F:fig1}
\end{figure}

La tabla \ref{T:tab1} es un ejemplo de como se elabora una tabla, además se puede usar el label para referenciar las tablas al igual que las imagenes.


%%%%% Tabla %%%%
\begin{table*}[!ht]
\caption{Comparación entre un transistor un BJT y un FET.}
\label{T:tab1}
\begin{center}
\begin{tabular}{| p{5.0cm} | p{3.5cm} | p{3.5cm} |}
\hline
\textbf{Características} & \textbf{BJT} & \textbf{FET}\\ 
\hline
Impedancia de entrada & Baja & Muy Alta \\ 
\hline
\end{tabular}
\end{center}
\end{table*}

Las ecuaciones se deben enumerar, equation hace esto automaticamente, se puede usar el label para referenciar las ecuaciones al igual que las imagenes. Como se muestra en la ecuación \ref{E:Ec1}.

\begin{equation}
\beta = \dfrac{I_C}{I_B}   
\label{E:Ec1}
\end{equation}


 Hay muchas formas de definir ecuaciones y argumentos en modelo matemático, en internet hay muchas referencias de como usar todas las variaciones.


\newpage


%%%%%%%%%%%%%%%%
\section{Diseño}
%%%%%%%%%%%%%%%%

\subsection{Primera parte.}

\subsubsection{Diseño de...}

Debe incluir los puntos de diseño propuestos en la hoja de experimentos y aquellos que el estudiante considere necesarios, con los cálculos, observaciones y recomendaciones pertinentes. Además debe incluir los circuitos esquemáticos y simulaciones que justifiquen el diseño propuesto. Algunos cálculos necesarios pueden incluirse en los Anexos si son muy grandes para esta sección.


\begin{figure}[H]
\centering
\includegraphics[width=0.9\textwidth]{./imagenes/fig2.png} 
\caption{Simulación de un...} 
\label{F:fig3}
\end{figure}






\newpage
%%%%%%%%%%%%%%%%%%%%%%%%%%
\section{Lista de equipos}
%%%%%%%%%%%%%%%%%%%%%%%%%%

La lista de equipos utilizados en el experimento se muestra en la Tabla \ref{T:equipos}.

\begin{table*}[!ht]
\caption{Lista de equipos}
\label{T:equipos}
\begin{center}
\begin{tabular}{| c | c | c | c | c | c | c |}
\hline
\textbf{Equipo} & \multicolumn{3}{| c |}{\textbf{Sesión 1}} & \multicolumn{3}{| c |}{\textbf{Sesión 2}} \\ 
\cline{2-7}
& \textbf{Marca} & \textbf{Modelo} & \textbf{Placa} & \textbf{Marca} & \textbf{Modelo} & \textbf{Placa} \\
\hline
Fuente DC &  &  &  & & &  \\
Osciloscopio digital &  &  &  & & &  \\
Generador de señales &  &  &  & & &  \\
Medidor Multifunción &  &  &  & & &  \\
Protoboard &  &  &  & & &  \\
\hline
\end{tabular}
\end{center}
\end{table*}
\newpage
%%%%%%%%%%%%%%%%%%%%%%%%%%%%%%
\section{Lista de componentes}
%%%%%%%%%%%%%%%%%%%%%%%%%%%%%%

La lista de componentes utilizados en el experimento se muestra en la Tabla \ref{T:componentes}.

\begin{table*} [!ht]
\caption{Lista de componentes}
\label{T:componentes}
\begin{center}
\begin{tabular}{| c | c | c | c | c | c |}
\hline
\textbf{Componente} & \textbf{Sigla} & \textbf{Valor nominal} & \textbf{Valor medido} & \textbf{Tolerancia} & \textbf{Potencia}\\ 
\hline
Resistor cerámico & R1 & \kiloohm[1] & &  & 0,25 W\\ 
Potenciómetro & Rp & \kiloohm[5] & &  & 2 W\\
Capacitor electrolítico & C11 & 150\microfarad & & - & - \\ 
Capacitor cerámico & C & 100\nanofarad & & - & - \\
Diodo & 1N4148 & & & & 500 mW\\
Transistor BJT & 2N2222A & & & & 1000 mW\\
\hline
\end{tabular}
\end{center}
\end{table*}


\newpage
%%%%%%%%%%%%%%%%%%%%%%%
\section{Procedimiento}
%%%%%%%%%%%%%%%%%%%%%%%


Este es un ejemplo de como debe verse un procedimiento, sea extenso y recuerde incluir todos los pasos necesarios para obtener los resultados que se esperan.



\begin{enumerate}
\item Encuentre la ganancia de corriente $\beta$ del transistor BJT, así como la corriente de saturación I$_{DSS}$ y la tensión de compuerta-fuente de apagado V$_p$ del transistor JFET.

\begin{itemize}
\item $\beta$ se calcula experimentalmente conectando el transistor en el modo emisor común y realizando mediciones de las corrientes de colector y de base, y con estas buscar su valor mediante la ecuación \ref{E:Ec1}.



\end{itemize}

Al finalizar esta sección se tendrán 4 capturas.

\item Determine la máxima excursión y el punto de operación de cada circuito. Compare los puntos de operación obtenidos con los diseñados.

La máxima exursión se da cuando la corriente de colector sea la mitad de su valor máximo.

\begin{table*} [!ht]
\caption{Datos máxima excursión}
\label{T:hola}
\begin{center}
\begin{tabular}{| c | c | c |}
\hline
\textbf{Diseño} & \textbf{Máxima excursión} & \textbf{Punto de operación} \\ 
\hline
Amplificador con FET &  &  \\ 
\hline
Amplificador con BJT &  &  \\ 
\hline
\end{tabular}
\end{center}
\end{table*}

\item Realice un barrido de frecuencias. Tome 10 medidas por década, desde 1 Hz hasta 100 kHz (50 puntos en total) y construya la 
gráfica de ganancia contra frecuencia. Determine las frecuencias de corte y el ancho de banda del circuito.

\begin{table*} [!ht]
\caption{Barrido de Frecuancias}
\label{T:hola}
\begin{center}
\begin{tabular}{| c | c |}
\hline
\textbf{Frecuencia} & \textbf{Ganancia} \\ 
\hline
1 Hz&  \\ 
\hline
... &  \\ 
\hline
100 Hz&  \\ 
\hline
... &  \\ 
\hline
1 kHz &  \\ 
\hline
... &  \\ 
\hline
10 kHz&  \\ 
\hline
... &  \\ 
\hline
100 kHz &  \\ 
\hline
... &  \\ 
\hline
\end{tabular}
\end{center}
\end{table*}

La frecuencia de corte se calcula cuando la ganancia es 0,707.
El ancho de banda es la diferencia entre las frecuencias de corte en alto y en bajo, se calcula con el osciloscopio.



\end{enumerate}

%%%%%%%%%%%%%%
% Bibliografía
%%%%%%%%%%%%%%

\newpage

El formato recomendado para la bibliografía es el APA. El siguiente es un ejemplo:

\begin{thebibliography}{}
\bibitem[Apuntes, 2008]{apuntes} \emph{http://www.apuntesdeelectronica.com/componentes/transistor-igbt.htm} consultado el 12/05/2013.
\bibitem[Boylestad, 1998]{Boylestad} Robert L. Boylestad, Louis Nashelsky (1998). \emph{Electronic Devices and Circuit Theory}. New Jersey:  Pearson Prentice Hall, 7th Edition.



\end{thebibliography}


\newpage

%%%%%%%%%%%%%%
% Anexos
%%%%%%%%%%%%%%

\appendix
\section{Anexos}

Incluye información de las hojas del fabricante de los componentes a utilizar o cualquier otra información que se considere necesaria.

La siguiente es una forma de como adjuntar páginas de PDFs con las hojas de fabricante, pero hay muchas formas de como adjuntar los anexos que se consideren necesarios.

%%%%% Transistor NTE123A y NTE159M
\subsection{Transistor 2N2222A}
\includepdf[pages={1-3}]{./Anexos/2N2222A.pdf}



%%%%%%%%%%%%%%
\end{document}
%%%%%%%%%%%%%%